% AY2025 制御工学 第2問〔II〕（転記）
% 出典: 03_制御工学/original/03_制御工学/2025/2025se.pdf
% 要約: 図2は Gc,Gp 直列の単位帰還系。補償器 Gc を変えた場合のステップ応答・
%       ゲイン線図を比較し, Gc の働きの違いを論じる。

\section*{第2問〔II〕}

制御システムに関する次の問い (1)〜(3) に答えよ. なお, 入力を $U(s)$, 出力を $Y(s)$ とし,
$G_c(s)$, $G_p(s)$ は伝達関数を表す. なお, 必要があれば $20\log_{10}\dfrac{3}{2}=3.52$ として
計算すること.

\begin{center}
\begin{tikzpicture}[>=Latex]
  \node (in)  at (0,0) {$U(s)$};
  \node (sum) [sum] at (1.3,0) {};
  \node (Gc)  [block] at (2.9,0) {$G_c(s)$};
  \node (Gp)  [block] at (4.9,0) {$G_p(s)$};
  \coordinate (nd) at (6.2,0);
  \node (out) at (7.2,0) {$Y(s)$};

  \draw[->] (in)  -- (sum);
  \draw[->] (sum) -- (Gc);
  \draw[->] (Gc)  -- (Gp);
  \draw[->] (Gp)  -- (out);
  \fill (nd) circle (1.4pt);
  \draw[->] (nd) -- (6.2,-1.3) -- (1.3,-1.3) -- (sum.south);

  \node at ($(sum.north west)+(0.02,0.10)$) {\scriptsize $+$};
  \node at ($(sum.south west)+(0.02,-0.10)$) {\scriptsize $-$};
  \node at (3.6,-1.7) {\small 図2};
\end{tikzpicture}
\end{center}

\begin{enumerate}[label=(\arabic*)]
  \item 図2に示す制御システムについて, 各伝達関数を $G_c(s)=\dfrac{2}{s+3}$,
        $G_p(s)=\dfrac{1}{s+1}$ とする場合のステップ応答 $y(t)$ を求めよ. また,
        一巡伝達関数に関するボード線図のうち, ゲイン線図を漸近線によって描け.

  \item 図2の制御システムにおいて $G_c(s)=\dfrac{s+10}{s}$, $G_p(s)=\dfrac{1}{s+1}$ とする
        場合のステップ応答 $y(t)$ を求めよ. また, 一巡伝達関数に関するボード線図のうち,
        ゲイン線図を漸近線によって描け. なお, 問い (1) で描いたゲイン線図との差異が明確に
        なるよう, 破線を用いて表すこと.

  \item 問い (1) ならびに問い (2) で得られたステップ応答とゲイン線図の差異に関する数値的な
        比較を行ったうえで, それぞれの制御システムにおける $G_c(s)$ の働きの違いについて
        詳細に説明せよ.
\end{enumerate}
